\section{Introduction}
% 1.4 改动
Large Language Models (LLMs) have demonstrated strong capabilities in natural language understanding and generation, but they are prone to hallucination and often lack up-to-date domain knowledge, which limits their reliability in knowledge-intensive applications~\cite{huang2025survey}. A common solution is Retrieval-Augmented Generation (RAG), which splits documents into passages, retrieves the top‑k relevant passages via dense vector search, and appends them to the prompt to ground the LLM’s response~\cite{gao2023retrieval, lewis2020retrieval}. This pipeline is simple, robust, and scalable, making it practical for real-world deployment. However, dense retrieval often treats seemingly unstructured text as flat content and overlooks the underlying structure that can be distilled from it~\cite{han2024retrieval}. This structured information is crucial for organizing and retrieving knowledge in massive, multi-source corpora, which limits the effectiveness of dense retrieval in realistic settings~\cite{peng2024graph}. Many queries require multi-step reasoning~\cite{plaat2025multi} and a deep understanding of the interconnections between entities and events across disparate sources. In these cases, RAG’s similarity-based top‑k retrieval emphasizes local paragraph relevance rather than knowledge structure, which can introduce redundant or unfocused context and reduce performance on multi-hop and relational questions~\cite{zhang2025survey}.

To bridge this structural gap, Graph-based RAG approaches have been proposed to incorporate Knowledge Graphs into the retrieval process. The standard workflow of Graph-based RAG~\cite{gao2023retrieval,zhang2025survey} typically involves two stages: Knowledge Graph Construction and Graph-Augmented Retrieval. In the construction phase, entities and their corresponding relations are extracted from raw corpus to build a structured graph~\cite{edge2024local, jimenez2024hipporag, gutierrez2025rag}, though some lightweight variants~\cite{zhuang2025linearrag} focus exclusively on entity indexing. During retrieval, these methods navigate the relational topology using graph algorithms such as Personalized PageRank (PPR)~\cite{yang2024efficient} to identify relevant evidence, which are then synthesized into textual context for the LLM.

% 1.4 改动
%不要出现Structure-centric
% 1 2递进关系，2解决1inefficiency但是损失accuracy，因为没有用relation
% 引出我们的，在考虑relation的基础上提升efficiency，同时保证high accuracy
However, existing solutions face a dilemma between accuracy and efficiency~\cite{han2025rag, xiang2025use}. \textbf{(1) Graph-based RAG methods can suffer from inefficient retrieval of information stored in knowledge graphs.} Methods such as HippoRAG~\cite{jimenez2024hipporag}, HippoRAG2~\cite{gutierrez2025rag} and ToG~\cite{sun2023think, ma2024think} preserve rich relational information for complex reasoning, but their reliance on graph traversal introduces substantial computational overhead and latency. Moreover, GraphRAG’s LLM-generated community summaries~\cite{edge2024local} can introduce semantic noise and secondary hallucinations, degrading retrieval precision. 
\textbf{(2) Some methods ignores edge information to simplify retrieval and improve efficiency, but at the cost of accuracy.} 
LinearRAG~\cite{zhuang2025linearrag} accelerates retrieval by constructing a relation-free hierarchical graph, termed Tri-Graph, using only lightweight entity extraction and semantic linking, effectively bypassing complex relational modeling. Similarly, $E^2$GraphRAG~\cite{zhao20252graphrag} streamlines the indexing process by utilizing SpaCy to construct bidirectional indexes between entities and passages to capture their many-to-many relations. While these approaches achieve significant speedups in indexing and retrieval, their reliance on simplified topologies weakens the explicit semantic dependencies required for precise multi-hop reasoning. Thus, a solution is still needed that maintains high accuracy while improving efficiency.


% P1
% To address these limitations, we propose [Method Name], a novel framework designed to synergize retrieval precision with computational efficiency. Our approach is grounded in the observation that triplets (subject, relation, object) serve as the fundamental atomic units of knowledge. Central to our framework is the introduction of \textbf{HyperNodes}, a higher-order representation where each triplet or its expanded relational path is treated as a unified node within a hypergraph structure. By shifting the initial search into this \textbf{Triple Space}, we effectively prune the vast search space in the embedding domain, filtering out irrelevant noise before document localization. This converts complex graph traversal into a streamlined process: first, pinpointing relevant knowledge triplets; and second, leveraging \textbf{Expanded-Path Retrieval} to dynamically aggregate these triplets into HyperNodes that represent explicit reasoning paths. This methodology preserves the structural integrity of knowledge graphs while maintaining the high efficiency of vector-based retrieval.

We propose \underline{\textbf{H}}yperNode \underline{\textbf{E}}xpansion and \underline{\textbf{L}}ogical \underline{\textbf{P}}ath-Guided Evidence Localization strategies for GraphRAG, named as \textbf{HELP}. It is a relation-aware retrieval framework that balances multi-hop accuracy with practical efficiency. On the one hand, we introduce HyperNode as a core unit to represent and chain knowledge triplets for multi-hop reasoning. In this framework, each retrieved knowledge triplet is instantiated as a HyperNode, which is then used to update the working graph, enabling iterative retrieval that chains HyperNodes together. This process naturally expands reasoning paths, transforming isolated facts into multi-hop relational chains.
On the other hand, we introduce Logical Path-Guided Evidence Localization, an efficient strategy that utilizes the expanded final HyperNodes as precise evidence to fetch the most relevant supporting passages from the corpus. By first constructing the relational path in the structured space, this approach converts the complex multi-hop question into a targeted retrieval task, ensuring that the returned passages are both semantically relevant and logically connected.
Extensive experiments show that our method can preserve the structural integrity of knowledge graphs while maintaining the high retrieval efficiency. Our contributions are summarized as follows:
%Our key idea is to treat knowledge triplets as first-class retrieval signals and to organize them into \textbf{HyperNodes}—higher-order units that bundle triples together with their expanded relational paths into unified retrieval entities. By explicitly representing relational chains rather than isolated text fragments, HyperNodes capture inter-fact dependencies and significantly strengthen multi-hop reasoning, while preserving strong performance on single-hop QA.

%To make this relation-aware retrieval efficient, we further introduce \textbf{Expanded-Path Retrieval}, a lightweight paradigm that operates directly in the embedding space. By shifting the initial search into the Triple Space, we effectively prune the vast search space in the embedding domain, filtering out irrelevant noise before document localization. This converts complex graph traversal into a streamlined process: first, pinpointing relevant knowledge triplets; and second, leveraging \textbf{Expanded-Path Retrieval} to dynamically aggregate these triplets into \textbf{HyperNodes} that represent explicit reasoning paths. This methodology preserves the structural integrity of knowledge graphs while maintaining the high efficiency of vector-based retrieval.



% 主打multihop qa，但simple qa也没有下降，引出HyperNode
%\textbf{1) HyperNode for Multi-Hop Reasoning:} 
\textbf{1) An Iterative HyperNode Retrieval Strategy for Reasoning Path Expansion.} We introduce the HyperNode, a higher-order retrieval unit that bundles triples together with their relational paths into a unified entity. By explicitly modeling knowledge paths rather than isolated fragments, HyperNode better captures inter-fact dependencies and substantially improves multi-hop reasoning, while maintaining strong performance on single-hop QA.
    
% 引出Expanded-Path
% \textbf{2) An Efficient Reasoning Path-Guided Evidence Localization Strategy.}
% We propose a lightweight Path-Guided Evidence Localization strategy. 
% By first constructing the reasoning path in structured triple space and then localizing evidence in the corpus, this strategy converts complex multi-hop questions into targeted retrieval, ensuring returned passages are both semantically relevant and logically connected—with high efficiency.

\textbf{2) An Efficient Reasoning Path-Guided Evidence Localization Strategy.} 
We propose a lightweight Path-Guided Evidence Localization strategy that leverages precomputed graph-text correlations. 
By mapping the HyperNodes constructed in the structured space directly to the source corpus, this strategy converts complex multi-hop questions into targeted retrieval tasks. 
% This mechanism bypasses computationally expensive graph traversals while ensuring that the returned passages are both semantically relevant and logically connected with superior efficiency.

% \textbf{3) Superior Performance and Adaptability.} Our framework achieves high accuracy with high computational efficiency. As a plug-and-play module, it can be seamlessly integrated into existing RAG pipelines without parameter updates, offering a robust and cost-effective enhancement across various Large Language Models.

% \textbf{3) Superior Performance and Parameter Robustness.} Our framework achieves a superior balance between multi-hop accuracy and computational efficiency. Crucially, our method demonstrates significant parameter robustness; empirical analysis confirms that HELP maintains consistent effectiveness across varying hyperparameter configurations, which offers a reliable and cost-effective enhancement for various Large Language Models.

\textbf{3) Superior Performance and Consistent Effectiveness.} Our framework achieves an excellent balance between multi-hop reasoning accuracy and computational efficiency. Crucially, our method demonstrates strong and consistent effectiveness across a variety of benchmarks under the same configuration. This consistent effectiveness renders our method a compelling, low-overhead enhancement across a wide array of Large Language Models.